% ID: SEL-D0010-BC-E06
% ESTADO: cotejada. Segunda pasada visual completa y borrador compilado revisados.
% Fuente: Selectividad/matematicas_aplicadas/EXTRA_TITULAR_MAT-APLIC-CCSS-JULIO-IMPRESO.pdf
% SHA-256: 548f60ee334516c7fa918dba3c28fb3ea5fe753cf76eb05402d7b8acb5a3b560
% Página: 2; bloque C; ejercicio original 6.

En una comunidad de vecinos, el 90\% de sus miembros tiene vehículo propio,
el 40\% hace uso del transporte público y un 3\% ni tiene vehículo propio
ni usa el transporte público. Se elige al azar un miembro de esa comunidad.
\begin{enumerate}[label=\alph*)]
  \item \textbf{(1 punto)} Calcule la probabilidad de que tenga vehículo
  propio o use el transporte público.
  \item \textbf{(0.5 puntos)} Calcule la probabilidad de que use el transporte
  público y no tenga vehículo propio.
  \item \textbf{(1 punto)} Calcule la probabilidad de que use el transporte
  público, sabiendo que no tiene vehículo propio.
\end{enumerate}
